%==============================================================================
\section{Phân tích hồi quy (Regression Analysis)}
\label{sec:regression-analysis}
%==============================================================================

\begin{dinhnghia}[title={Phân tích hồi quy}]
Kỹ thuật thống kê dự đoán \textbf{giá trị số liên tục} dựa trên quan hệ giữa biến độc lập và biến phụ thuộc.
Mục tiêu: vẽ đường/đường cong khớp dữ liệu tốt nhất và ước lượng ảnh hưởng của biến này lên biến kia.
\end{dinhnghia}

\begin{ynghia}[title={Vai trò trong học máy}]
Hồi quy là \textbf{học có giám sát}: mô hình học từ feature (độc lập) và target (phụ thuộc, số liên tục) để dự đoán trên dữ liệu mới.
Ứng dụng: dự báo, phân tích dự đoán, tối ưu kinh doanh, kinh tế, tài chính.
\end{ynghia}

\subsection{Thuật ngữ}

\begin{phuongphap}[title={Thuật ngữ thường gặp}]
\begin{itemize}
  \item \textbf{Biến độc lập} (predictor, feature): dùng để dự đoán.
  \item \textbf{Biến phụ thuộc} (target, response): giá trị cần dự đoán.
  \item \textbf{Đường hồi quy}: đường/đường cong khớp điểm dữ liệu.
  \item \textbf{Overfitting}: khớp tốt tập huấn luyện, kém tập kiểm tra (phương sai cao).
  \item \textbf{Underfitting}: mô hình kém cả trên tập huấn luyện (độ lệch cao).
  \item \textbf{Outlier}: điểm cực cao/thấp, lệch pattern chung.
  \item \textbf{Đa cộng tuyến} (multicollinearity): feature phụ thuộc lẫn nhau.
\end{itemize}
\end{phuongphap}

\subsection{Cách hồi quy hoạt động}

\begin{phuongphap}[title={Quy trình}]
Thuật toán hồi quy được huấn luyện trên tập có nhãn (feature + target).
Trong huấn luyện, mô hình học quan hệ giữa predictor và target; sau đó dự đoán giá trị mới từ quan hệ đã học.
\end{phuongphap}

\subsection{Các loại hồi quy trong học máy}

Phân loại thường dựa trên: số biến độc lập, kiểu biến phụ thuộc, hình dạng đường hồi quy.

\begin{phuongphap}[title={Các loại phổ biến}]
\begin{enumerate}
  \item \textbf{Hồi quy tuyến tính} --- quan hệ tuyến tính; gồm đơn biến và đa biến.
  \item \textbf{Hồi quy logistic} --- xác suất sự kiện; phân loại nhị phân (sigmoid).
  \item \textbf{Hồi quy đa thức} --- quan hệ bậc $n$; bắt phi tuyến, nhạy outlier.
  \item \textbf{Lasso} (L1) --- chuẩn hóa, giảm overfitting; dữ liệu chiều cao.
  \item \textbf{Ridge} (L2) --- chuẩn hóa; xử lý đa cộng tuyến, nhiều tham số.
  \item \textbf{Cây quyết định hồi quy} --- chia tập con, giảm MSE từng nút.
  \item \textbf{Random Forest hồi quy} --- tập hợp nhiều cây (bagging).
  \item \textbf{SVR} --- máy vectơ hỗ trợ cho hồi quy; chịu outlier, phi tuyến.
\end{enumerate}
\end{phuongphap}

\begin{dinhnghia}[title={Hồi quy tuyến tính đơn}]
Một biến độc lập $X$ dự đoán $Y$:
\[
Y = mX + b
\]
($Y$ phụ thuộc, $X$ độc lập, $m$ độ dốc, $b$ hệ số chặn.)
\end{dinhnghia}

\begin{luuy}[title={Lỗi trong nguồn gốc}]
Bài gốc Tutorialspoint đôi chỗ ghi nhầm $X$/$Y$ là biến phụ thuộc/độc lập; phần dưới dùng ký hiệu chuẩn.
\end{luuy}

\subsection{Mô hình hồi quy đơn và đa biến}

\begin{itemize}
  \item \textbf{Mô hình đơn giản}: một feature duy nhất.
  \item \textbf{Mô hình đa biến}: nhiều feature.
\end{itemize}

\subsection{Chọn mô hình hồi quy}

\begin{phuongphap}[title={Tiêu chí}]
So sánh MSE, MAE, $R^2$; độ phức tạp; khả năng giải thích.
Chọn mô hình có metric tốt, đơn giản và dễ diễn giải khi có thể.
\end{phuongphap}

\subsection{Metric đánh giá}

\begin{tinhchat}[title={Các chỉ số}]
\begin{itemize}
  \item \textbf{MAE}: trung bình $|y - \hat{y}|$.
  \item \textbf{MSE}: trung bình $(y - \hat{y})^2$.
  \item \textbf{Median absolute error}: trung vị $|y - \hat{y}|$.
  \item \textbf{RMSE}: $\sqrt{\mathrm{MSE}}$.
  \item \textbf{$R^2$}: tối đa 1.0; có thể âm nếu mô hình rất kém.
  \item \textbf{MAPE}: MAE dạng phần trăm.
\end{itemize}
\end{tinhchat}

\subsection{Ứng dụng}

\begin{ynghia}[title={Ứng dụng điển hình}]
\begin{itemize}
  \item \textbf{Dự báo}: GDP, giá dầu, chuỗi thời gian.
  \item \textbf{Tối ưu}: giờ cao điểm khách hàng.
  \item \textbf{Hiệu chỉnh quyết định}: đánh giá và sửa chính sách.
  \item \textbf{Kinh tế / tài chính}: cung--cầu, rủi ro danh mục.
\end{itemize}
\end{ynghia}

\subsection{Xây dựng regressor trong Python}

\begin{phuongphap}[title={Các bước (scikit-learn)}]
\begin{enumerate}
  \item Import thư viện.
  \item Nạp dữ liệu.
  \item Chia train/test.
  \item Huấn luyện và dự đoán.
  \item Vẽ đồ thị.
  \item Tính metric.
\end{enumerate}
\end{phuongphap}

\begin{vidu}[title={Hồi quy tuyến tính với dữ liệu tổng hợp}]
Thay đường dẫn \texttt{C:\textbackslash linear.txt} và lỗi \texttt{num\_training} trong bài gốc bằng \texttt{make\_regression}:

\begin{lstlisting}
import numpy as np
from sklearn import linear_model
import sklearn.metrics as sm
import matplotlib.pyplot as plt
from sklearn.datasets import make_regression

X, y = make_regression(n_samples=100, n_features=1, noise=15, random_state=42)

training_samples = int(0.6 * len(X))
X_train, y_train = X[:training_samples], y[:training_samples]
X_test, y_test = X[training_samples:], y[training_samples:]

reg_linear = linear_model.LinearRegression()
reg_linear.fit(X_train, y_train)
y_test_pred = reg_linear.predict(X_test)

plt.scatter(X_test, y_test, color='red')
plt.plot(X_test, y_test_pred, color='black', linewidth=2)
plt.show()

print("Regressor model performance:")
print("MAE =", round(sm.mean_absolute_error(y_test, y_test_pred), 2))
print("MSE =", round(sm.mean_squared_error(y_test, y_test_pred), 2))
print("R2 =", round(sm.r2_score(y_test, y_test_pred), 2))
\end{lstlisting}
\end{vidu}
